\section{설계 목적과 연구 기반}
\label{sec:introduction}
뇌--컴퓨터 인터페이스(BCI)는 신호 획득, 신호 변환, 상태 추정, 사건 생성과 피드백을 연결한다.
하나의 실험은 전극 전압과 광학 측정, 입력 사건과 모델 상태, 여러 장치의 시계를 함께 다룬다.
본 설계는 이 관계들을 언어의 형식과 실행 의미에 포함한다. 프로그램은 필요한 측정량과
시간 관계를 기술하고, 컴파일러는 선택한 하드웨어 구성과 연산자의 계약을 결합하여 실행 계획을 생성한다.

BCI2000은 신호 획득, 신호 처리, 응용 기능을 모듈로 구성하고 실험 조건과 상태를 공유하는
구조를 제시하였다\cite{bci2000}. BCILAB은 신호 처리와 학습 절차를 조합하고 방법을 평가하는
프레임워크를 구성하였다\cite{bcilab}. MOABB는 데이터셋, 전처리와 평가 인터페이스를 통합하여
BCI 알고리즘 비교의 재현성을 다룬다\cite{moabb}. 본 논문은 이 세 작업에서 사용하는
획득 경계, 처리 그래프, 평가 단위를 명시적 언어 요소로 연결한다.

형식적 기반은 스트림 의미론, 동기 데이터흐름, 소유권에 둔다. Kahn의 프로세스 네트워크는
통신 이력을 접두 순서와 함수 합성으로 해석한다\cite{kahn}. Synchronous Data Flow(SDF)는
연산자의 생산ㆍ소비 토큰 수를 이용하여 다중 표본률 실행 계획을 구성한다\cite{sdf}.
Lustre는 논리 시계와 지연을 반응형 프로그램의 구조에 포함한다\cite{lustre}.
RustBelt는 소유권과 차용의 의미 및 라이브러리 경계의 검증 조건을 제시한다\cite{rustbelt}.
본 설계는 스트림에 측정 시각과 가용 시각을 부착하고, 각 가변 자원에 소유 actor를 배정하며,
시계ㆍ학습ㆍ하드웨어 계약을 그래프의 합성 조건으로 사용한다.

설계의 주요 산출물은 다섯 가지이다. 첫째, 물리 차원과 명목 채널 의미를 포함한 스트림 형식이다.
둘째, 상태 전이와 가용 시각 인과성을 결합한 코어 의미론이다. 셋째, 구성 공간과 기능 집합을
사용하는 하드웨어 바인딩이다. 넷째, 시계 변환ㆍ창 완결ㆍ버퍼 용량의 계산 규칙이다.
다섯째, 모델 계보와 실행 사건을 재생 가능한 기록으로 묶는 Rust 실행 구조이다.

\subsection{정의, 근거, 계산의 구성}
논문은 정의와 계약으로 언어의 동작을 지정하고, 인용 문헌으로 장치 규격과 선행 연구의
근거를 연결한다. 명제는 지정한 코어 의미론에서 합성 성질을 증명한다. 수치 예제는
장치 규격과 명시한 실행 매개변수를 식에 대입하여 처리량ㆍ지연ㆍ용량을 계산한다.
각 하드웨어 프로필은 원문의 기관, 제목, 버전 또는 열람일을 함께 가진다.
2026년 10월 31일의 문서 판본을 참고문헌에 기록한다.

\begin{figure}[tb]
\centering
\begin{tikzpicture}[>=Latex,font=\small,node distance=15mm]
\node (source) at (0,1.4) {프로그램과 프로필};
\node (ir) at (4.6,1.4) {형식화 중간 표현};
\node (exec) at (9.2,1.4) {Rust 실행 계획};
\node (contract) at (0,-0.3) {측정ㆍ시간ㆍ기능 계약};
\node (sem) at (4.6,-0.3) {스트림 의미와 불변식};
\node (trace) at (9.2,-0.3) {값ㆍ사건ㆍ계보 기록};
\draw[->] (source)--node[above]{구체화}(ir);
\draw[->] (ir)--node[above]{계획 생성}(exec);
\draw[->] (source)--(contract);
\draw[->] (ir)--(sem);
\draw[->] (exec)--(trace);
\draw[->] (contract)--node[below]{합성 조건}(sem);
\draw[->] (sem)--node[below]{실행 관계}(trace);
\end{tikzpicture}
\caption{언어 구체화와 실행 기록의 대응. 위쪽 경로는 컴파일 과정을, 아래쪽 경로는 각 산출물에 붙는 의미와 계약을 나타낸다.}
\label{fig:architecture}
\end{figure}

\subsection{언어의 조직}
표면 언어는 장치 프로필 선택, 채널 스키마, 신호 연산, 학습 프로토콜, 시간 계약을 선언한다.
컴파일러는 유리수 표본률, 정수 창 길이, 유한 채널 집합과 명목 식별자를 정규화한다.
구성 조회로 얻는 값은 런타임 기술자와 연결한다. 이후 형식화 중간 표현은 원시 연산자,
직렬ㆍ병렬 합성, 지연, 피드백으로 구성한다. 이 조직은 새로운 장치와 알고리즘을
프로필과 원시 연산자 계약으로 추가하는 확장 지점을 제공한다.

측정값의 물리 단위, 장치 시각, 모델의 데이터 계보는 각각 독립적인 정보이다.
언어는 이 정보를 동일한 처리 그래프 안에 배치하고 변환마다 입력ㆍ출력 관계를 기록한다.
예를 들어 전압 단위 변환은 수치 배율을, 재기준화는 채널 스키마를,
시계 정렬은 시간 뷰를, 모델 갱신은 매개변수 버전과 가용 시각을 변경한다.
